% Part: first-order-logic
% Chapter: introduction
% Section: formulas

\documentclass[../../../include/open-logic-section]{subfiles}

\begin{document}

\olfileid{fol}{int}{fml}

% SEGMENT OLP-0142-S01
\section{\usetoken{P}{formula}}

முதல்தர தருக்கத்தின் !!{sentence}s ஐக் கடுமையாகக் குறிப்பிடவும்,
!!{variable}s பயன்பாட்டால் எழும் சிக்கல்களைக் கையாளவும் நாம் பின்வரும்
அணுகுமுறையைப் பயன்படுத்துவோம். முதலில் \emph{வேறொரு} கோவைகளின் கணத்தை,
அதாவது !!{formula}s ஐ வரையறுக்கிறோம். அதைச் செய்ததும், அவற்றில்
!!{variable}s வகிக்கும் பங்கைக் கருதலாம்; ``கட்டற்ற'' மற்றும் ``கட்டுண்ட''
!!{variable}s என்ற வேறு சில கருத்துகளின் அடிப்படையில் !!a{sentence} என்பதை
வரையறுக்கலாம் (அதாவது, கட்டற்ற !!{variable}s இல்லாத !!a{formula}).
``!!{sentence}'' என்பதன் வரையறையை இது எளிதாக்குவதால் மட்டுமல்லாமல்,
நிறைவுறுத்தல் என்ற பொருண்மையியல் கருத்தை நாம் வரையறுக்கும் முறைக்கு இது
இன்றியமையாதது என்பதாலும் இவ்வாறு செய்கிறோம்.

% SEGMENT OLP-0142-S02
ஒரே ஒரு !!{predicate}~$\Obj P$, ஒரே ஒரு !!{constant}~$\Obj a$ மற்றும்
$\lnot$, $\land$, $\lexists$ ஆகிய தருக்கக் குறிகளை மட்டும் கொண்ட ஓர் எளிய
முதல்தர மொழிக்காக ``!!{formula}'' என்பதை வரையறுப்போம். நமது முழு வரையறைகள்
இதைவிட மிகவும் பொதுவானவை: முடிவிலி எண்ணிக்கையிலான !!{predicate}s மற்றும்
!!{constant}s ஐ அனுமதிப்போம். உண்மையில், ``சொற்கள்'' உருவாகுமாறு
!!{constant}s மற்றும் !!{variable}s உடன் சேர்க்கக்கூடிய !!{function}s ஐயும்
கருதுவோம். இப்போதைக்கு~$\Obj a$ மற்றும் மாறிகள் மட்டுமே நமது சொற்கள்.
!!{variable}s மட்டும் முடிவிலி எண்ணிக்கையில் தேவை. $\Obj v_0$, $\Obj v_1$,
\dots{} என்ற குறிகளை அதிகாரபூர்வமாக மாறிகளாகப் பயன்படுத்துவோம்.

% SEGMENT OLP-0142-S03
\begin{defn}
\emph{!!{formula}s} கணம்~$\Frm$ பின்வருமாறு வரையறுக்கப்படுகிறது:
\begin{enumerate}
\item\ollabel{fmls-atom} $\Atom{\Obj P}{\Obj a}$ மற்றும் $\Atom{\Obj
  P}{\Obj v_i}$ ஆகியவை !!{formula}s ($i \in \Nat$).

\tagitem{prvNot}{\ollabel{fmls-not}$!A$ ஓர் !!a{formula} எனில், $\lnot !A$
  ஒரு !!{formula}.}{}

\tagitem{prvAnd}{$!A$ மற்றும் $!B$ ஆகியவை !!{formula}s எனில், $(!A \land
  !B)$ ஓர் !!a{formula}.}{}

\tagitem{prvEx}{\ollabel{fmls-ex}$!A$ ஓர் !!a{formula} ஆகவும், $x$ ஓர்
  !!a{variable} ஆகவும் இருந்தால், $\lexists[x][!A]$ ஓர் !!a{formula}.}{}

\tagitem{limitClause}{\ollabel{fmls-limit}வேறெதுவும் !!a{formula} அல்ல.}{}
\end{enumerate}
\end{defn}

% SEGMENT OLP-0142-S04
எந்த $i \in \Nat$ க்கும் $\Atom{\Obj P}{\Obj a}$ மற்றும்
$\Atom{\Obj P}{\Obj v_i}$ ஆகியவை !!{formula}s என்று~\olref{fmls-atom}
கூறுகிறது. இவை \emph{அணு} !!{formula}s எனப்படுகின்றன. இவை நமக்குத்
தொடக்கப் புள்ளியைத் தருகின்றன. மற்ற கூறுகள், நாம் ஏற்கெனவே உருவாக்கிய
கோவைகளிலிருந்து புதிய !!{formula}s ஐ உருவாக்கும் வழிகளைத் தருகின்றன.
எடுத்துக்காட்டாக, $\Atom{\Obj P}{\Obj v_2}$ ஏற்கெனவே !!a{formula} என்பதால்,
$\lnot \Atom{\Obj P}{\Obj v_2}$ உம் !!a{formula} என்று~\olref{fmls-not}
படி கிடைக்கிறது; இது~\olref{fmls-atom} படி அடிப்படை ஆகும்.
பின்னர்
$\lexists[\Obj v_2][\lnot \Atom{\Obj P}{\Obj v_2}]$ வேறொரு !!{formula}
என்று~\olref{fmls-ex} படி கிடைக்கிறது; இவ்வாறே தொடரலாம். இவ்வழியில் உருவாக்க
முடியும் குறித்தொடர்கள் \emph{மட்டுமே} !!{formula}s ஆகக் கணக்கிடப்படும்
என்று~\olref{fmls-limit} கூறுகிறது. குறிப்பாக,
$\lexists[\Obj v_0][\Atom{\Obj P}{\Obj a}]$ மற்றும் $\lexists[\Obj
  v_0][\lexists[\Obj v_0][\Atom{\Obj P}{\Obj a}]]$ ஆகியவை
!!{formula}s ஆக \emph{உண்மையிலேயே} கணக்கிடப்படும். ஆனால்
$(\lnot \Atom{\Obj P}{\Obj a})$ மிகையாக வெளிப்புற அடைப்புக்குறிகளைக்
கொண்டிருப்பதால் அவ்வாறு கணக்கிடப்படாது.

% SEGMENT OLP-0142-S05
!!{formula}s ஐ இவ்வாறு வரையறுப்பது \emph{தொகுத்தறிதல் வழி வரையறை}
எனப்படுகிறது. \emph{கட்டமைப்புத் தொகுத்தறிதல்} எனப்படும் தொகுத்தறிதல் நிறுவல்
வடிவத்தைப் பயன்படுத்தி !!{formula}s பற்றியவற்றை நிறுவ இது உதவுகிறது.
இவை~\olref[mth][ind][idf]{sec} மற்றும் \olref[mth][ind][sti]{sec} ஆகியவற்றில்
பொதுவாக விவாதிக்கப்பட்டுள்ளன; பிந்தைய நிறுவல்களை ஆழமாகப் படிப்பதற்கு முன்
அவற்றை மீள்பார்வையிட வேண்டும். அடிப்படையில், எல்லா !!{formula}s க்கும் ஒன்று
மெய் என்று நிறுவ விரும்பினால், முதலில் அது அணு !!{formula}s க்கு மெய் எனக்
காட்டுகிறோம். பின்னர் எந்த !!{formula}~$!A$ க்கும் (மற்றும்~$!B$ க்கும்) அது
மெய் \emph{எனில்}, $\lnot !A$, $(!A \land !B)$ மற்றும் $\lexists[x][!A]$
ஆகியவற்றுக்கும் அது \emph{மெய்} எனக் காட்டுகிறோம். எடுத்துக்காட்டாக, இது
$\lexists[\Obj v_2][\lnot \Atom{\Obj P}{\Obj v_2}]$ க்கு மெய் என்று
நிறுவுகிறது: முதல் பகுதியிலிருந்து அணு !!{formula}~$\Atom{\Obj P}{\Obj v_2}$
க்கு அது மெய் என அறிவோம். பின்னர் இரண்டாம் பகுதியால்
$\lnot \Atom{\Obj P}{\Obj v_2}$ க்கு அது மெய் என்றும், மீண்டும் அதனால்
$\lexists[\Obj v_2][\lnot \Atom{\Obj P}{\Obj v_2}]$ க்கே அது மெய் என்றும்
கிடைக்கிறது. எல்லா !!{formula}s உம் அணு !!{formula}s இலிருந்து தொகுத்தறிதல்
வழியாக உருவாக்கப்படுவதால், அவற்றில் எதற்கும் இது செயல்படும்.

\end{document}
